% !PRINTED 316
\weberpage{316}

\begin{center}
\S.\ 104.

\gesperrt{Determination of an upper limit for the roots.}
\end{center}

We have already, in the third section, made use of the theorem that
one can always find a positive number which is greater than the
absolute values of all the roots of a given equation. But now the
point is to give a determination of such an upper limit that is as
simple and at the same time as accurate as possible.

We consider at first only the real roots, and thus seek a positive
number which is greater than the greatest positive root of a given
equation $f(x)=0$.

Such a determination is given to us by the theorem of \gesperrt{Budan}.

If we take the coefficient of the highest power of $x$ in $f(x)$
equal to 1 (or at least positive), then the positive number
$\alpha$ can always be taken so large that
\weq{1}{f(\alpha),\ f'(\alpha),\ f''(\alpha)\ldots f^{(n)}(\alpha)}
all become positive. Then in the sequence of functions
\[
f(x),\ f'(x),\ f''(x)\ldots f^{(n)}(x)
\]
no change of sign is lost any more between $x=\alpha$ and $x=\infty$, and
hence, by Theorem I, no root of $f(x)$ can lie between
$\alpha$ and $\infty$.

If one wishes to obtain from the same theorem a lower limit for the
negative roots, one takes the positive number $\beta$ so that
\weq{2}{(-1)^n f(-\beta),\ (-1)^{n-1} f'(-\beta),\ \ldots f^{(n)}(-\beta)}
all become positive; then the equation $f(x)=0$ certainly has no
negative root below $-\beta$.

This determination of the limits goes back already to \gesperrt{Newton}.

Another determination of an upper limit, which is simpler for
computation, has been given by \gesperrt{Laguerre}\footnote{\gesperrt{Laguerre},
Nouvelles Annales d.\ math., 2.\ ser., t.\ XIX, 1880. Journal de math.,
3.\ ser., t.\ IX, 1883.}.

Instead of the derivatives $f(x), f'(x), f''(x)\ldots$ he uses the
functions $f_0(x), f_1(x), f_2(x), \ldots f_{n-1}(x)$, which have already
%CONT>>
